\documentclass[11pt]{article}
\usepackage{amsmath, amssymb}
\usepackage{booktabs}
\usepackage[margin=1in]{geometry}
\usepackage{hyperref}

\title{Reward Function and State Space\\
\large Octopus CXL Memory Pooling --- RL Redesign}
\author{Pulak Mehrotra}
\date{}

\begin{document}
\maketitle

% -----------------------------------------------------------------------
\section*{Notation}

\begin{tabular}{ll}
$m$                    & number of MPDs in the pod \\
$[m]$                  & set $\{1, \ldots, m\}$ \\
$N(i)$                 & MPDs reachable from host $i$ \\
$H(j)$                 & hosts connected to MPD $j$ \\
$\deg(h)$              & degree of host $h$, i.e.\ $|N(h)|$ \\
$c_j(t)$               & current allocated load on MPD $j$ at time $t$ (GB) \\
$\ell_h(t)$            & current CXL load sourced by host $h$ at time $t$ (GB) \\
$\mathcal{V}_j(t)$     & VMs currently allocated on MPD $j$ at time $t$ \\
$\mathrm{mem}(v)$      & allocated memory of VM $v$ (GB) \\
$\mathrm{end}(v)$      & predicted departure time of VM $v$;
                         computed from trace during training,
                         and from a lifetime prediction model
                         at deployment \cite{pond2023} \\
$x$                    & memory request of the arriving VM (GB) \\
$\mathbf{w}$           & allocation weight vector,
                         $w_j \geq 0$,
                         $\sum_{j \in N(i)} w_j = 1$,
                         $w_j = 0$ if $j \notin N(i)$ \\
$W$                    & lookahead window (timesteps; target $150$--$300$) \\
$D_\mathrm{pod}$       & total pod DRAM capacity (GB); used for normalisation \\
$\lambda$              & Reward-B global-term weight ($\lambda < 1$;
                         start $0.1$--$0.3$) \\
$d_\mathrm{max}$       & maximum host degree in the topology \\
\end{tabular}

% -----------------------------------------------------------------------
\section{Reward Function}

\subsection{Departure Relief}

Define the time-weighted departure relief on MPD $j$ over the
lookahead window $[t,\; t{+}W]$:

\begin{equation}
    D_j(t, W)
    \;=\;
    \sum_{\substack{v \,\in\, \mathcal{V}_j(t) \\
                    \mathrm{end}(v) \;\leq\; t + W}}
    \mathrm{mem}(v)
    \cdot
    \left(1 - \frac{\mathrm{end}(v) - t}{W}\right)
    \label{eq:departure}
\end{equation}

Each VM's contribution is weighted by how soon it departs: a VM
departing at $t{+}1$ receives weight $\approx 1$ (near-certain
relief); a VM departing at $t{+}W$ receives weight $\approx 0$.
VMs departing after $t{+}W$ do not contribute to $D_j$ and
instead constitute the sticky load (Section~\ref{sec:state}).

\subsection{Projected Load}

The projected load on MPD $j$ before placement is:

\begin{equation}
    \hat{c}_j(t)
    \;=\;
    \frac{c_j(t) \;-\; D_j(t, W)}{D_\mathrm{pod}}
    \label{eq:chat}
\end{equation}

After the agent applies allocation $\mathbf{w}$, the projected load
on each reachable MPD $j \in N(i)$ is updated:

\begin{equation}
    \hat{c}_j^{\,+}(t)
    \;=\;
    \hat{c}_j(t)
    \;+\;
    \frac{w_j \cdot x}{D_\mathrm{pod}}
    \label{eq:chat-post}
\end{equation}

\subsection{Reward}

\paragraph{Reward A: local projected-peak baseline.}
\begin{equation}
\boxed{
    R_A(t)
    \;=\;
    -\,\max_{j \,\in\, N(i)}\, \hat{c}_j^{\,+}(t)
}
\label{eq:reward-a}
\end{equation}

This is the default baseline reward.  It penalises the worst projected
load among the MPDs the current host can directly influence, keeping the
signal local and easy to interpret.

\paragraph{Reward B: local projected peak plus coarse global term.}
\begin{equation}
\boxed{
    R_B(t)
    \;=\;
    -\,\max_{j \,\in\, N(i)}\, \hat{c}_j^{\,+}(t)
    \;-\;
    \lambda \cdot \max_{j \,\notin\, N(i)}\, \hat{c}_j(t)
}
\label{eq:reward-b}
\end{equation}

Reward B is an ablation variant that adds a coarse global-stress term.
The first term is identical to Reward A.  The second term, down-weighted
by $\lambda$, penalises the projected peak among unreachable MPDs.  Its
purpose is to test whether coarse pod-level stress helps value
estimation beyond the local baseline.

All values are normalised by $D_\mathrm{pod}$.  Reward A lies in
$(-1, 0]$.  Reward B lies in $(-(1+\lambda), 0]$.  This keeps both
reward variants on a stable scale across traces and pod sizes.

% -----------------------------------------------------------------------
\section{State Space}
\label{sec:state}

\subsection{Design Principle}

The state is shared across Reward A and Reward B.  Its main job is to
describe, for each reachable MPD, both
\emph{(i)} how much load is currently there and how much of that load is
likely to drain soon, and
\emph{(ii)} how contested that MPD is by nearby hosts.
This keeps the baseline design local and interpretable while exposing
the neighborhood structure that greedy does not model explicitly.

\subsection{Per-MPD Features}

The following features are computed for each MPD slot
$k = 1, \ldots, d_\mathrm{max}$, zero-padded for slots
$k > |N(i)|$:

\begin{align}
    \phi_k^{(1)} &\;=\;
        \frac{c_{j_k}(t)}{D_\mathrm{pod}}
        && \text{current allocated load on MPD } j_k
        \label{eq:phi1} \\[6pt]
    \phi_k^{(2)} &\;=\;
        \frac{D_{j_k}(t,\,W)}{D_\mathrm{pod}}
        && \text{time-weighted departures within } W \text{ steps}
        \label{eq:phi2} \\[6pt]
    \phi_k^{(3)} &\;=\;
        \frac{S_{j_k}(t)}{D_\mathrm{pod}}
        && \text{sticky load (departs after } W \text{ steps)}
        \label{eq:phi3} \\[6pt]
    \phi_k^{(4)} &\;=\;
        \mathbf{1}\bigl[j_k \in N(i)\bigr]
        && \text{connectivity mask / action validity}
        \label{eq:phi4} \\[6pt]
    \phi_k^{(5)} &\;=\;
        \frac{P_{j_k}(t)}{D_\mathrm{pod}}
        && \text{neighborhood pressure on MPD } j_k
        \label{eq:phi5} \\[6pt]
    \phi_k^{(6)} &\;=\;
        Q_{j_k}
        && \text{neighbor scarcity around MPD } j_k
        \label{eq:phi6}
\end{align}

where the sticky load is defined as:

\begin{equation}
    S_j(t)
    \;=\;
    \sum_{\substack{v \,\in\, \mathcal{V}_j(t) \\
                    \mathrm{end}(v) \;>\; t + W}}
    \mathrm{mem}(v)
    \label{eq:sticky}
\end{equation}

and the two neighborhood-contention features are defined as:

\begin{align}
    P_j(t)
    &\;=\;
    \sum_{h \in H(j)} \ell_h(t)
    && \text{total current host pressure from hosts attached to } j
    \label{eq:pressure} \\[6pt]
    Q_j
    &\;=\;
    \frac{1}{|H(j)|}
    \sum_{h \in H(j)}
    \frac{1}{\deg(h)}
    && \text{average scarcity of alternatives among hosts attached to } j
    \label{eq:scarcity}
\end{align}

Note that $\phi_k^{(1)} = \phi_k^{(2)} + \phi_k^{(3)}$ up to
normalisation: current load decomposes exactly into draining and
sticky components.  The agent is therefore given the full
decomposition of $\hat{c}_j$ directly rather than having to
infer it from a single load snapshot.  The added
$\phi^{(5)}$ and $\phi^{(6)}$ terms provide a simple 2-hop summary of
which reachable MPDs are likely to be needed by neighboring hosts.

\subsection{Global Context Features}

\begin{align}
    \psi^{(1)} &\;=\;
        \frac{\displaystyle\max_{j \in [m]}\, c_j(t)}{D_\mathrm{pod}}
        && \text{global peak load (used only by Reward B)}
        \label{eq:psi1} \\[6pt]
    \psi^{(2)} &\;=\;
        \frac{x}{D_\mathrm{pod}}
        && \text{arriving VM size}
        \label{eq:psi2}
\end{align}

\subsection{Full Observation Vector}

\begin{equation}
    s_t
    \;=\;
    \Bigl[
        \phi_1^{(1:6)},\;
        \phi_2^{(1:6)},\;
        \ldots,\;
        \phi_{d_\mathrm{max}}^{(1:6)},\;
        \psi^{(1)},\;
        \psi^{(2)}
    \Bigr]
    \;\in\; \mathbb{R}^{\,6 d_\mathrm{max} \,+\, 2}
    \label{eq:obs}
\end{equation}

For the AG16x6 topology ($d_\mathrm{max} = 8$) this yields a
$\mathbf{50}$-dimensional observation vector.

\subsection{State--Reward Alignment}

\begin{table}[h]
\centering
\caption{Shared state features used by Reward A and Reward B.}
\label{tab:alignment}
\begin{tabular}{lll}
\toprule
Feature & Reward quantity & Role \\
\midrule
$\phi^{(1)}$   & $c_j / D_\mathrm{pod}$
               & Starting point for $\hat{c}_j$ \\
$\phi^{(2)}$   & $D_j(t,W) / D_\mathrm{pod}$
               & Draining component of $\hat{c}_j$ \\
$\phi^{(3)}$   & $S_j(t) / D_\mathrm{pod}$
               & Permanent component of $\hat{c}_j$ \\
$\phi^{(4)}$   & Action mask
               & Enforces topology constraint \\
$\phi^{(5)}$   & $P_j(t) / D_\mathrm{pod}$
               & Neighborhood contention around MPD $j$ \\
$\phi^{(6)}$   & $Q_j$
               & Scarcity of alternatives for nearby hosts \\
$\psi^{(1)}$   & $\max_j c_j / D_\mathrm{pod}$
               & Coarse global context for Reward B only \\
$\psi^{(2)}$   & $x / D_\mathrm{pod}$
               & Post-placement update magnitude \\
\bottomrule
\end{tabular}
\end{table}

\subsection{Comparison with v0}

\begin{table}[h]
\centering
\caption{State space comparison.}
\label{tab:comparison}
\begin{tabular}{lll}
\toprule
Feature & v0 & This work \\
\midrule
Current MPD load          & \checkmark & \checkmark \\
Connectivity mask         & \checkmark & \checkmark \\
Arriving VM size          & \checkmark & \checkmark \\
Global peak load          & \checkmark & \checkmark \\
Departure relief $D_j(W)$ & $\times$   & \checkmark \\
Sticky load $S_j$         & $\times$   & \checkmark \\
Neighborhood pressure $P_j$ & $\times$ & \checkmark \\
Neighbor scarcity $Q_j$   & $\times$   & \checkmark \\
\midrule
Dimensions (AG16x6)       & 20         & 50         \\
\bottomrule
\end{tabular}
\end{table}

Relative to v0, the added dimensions come from two sources:
future-pressure decomposition $(D_j, S_j)$ and neighborhood-contention
signals $(P_j, Q_j)$.  Together they tell the policy not only what load
may drain soon, but also which reachable MPDs are strategically valuable
because they serve a more constrained neighborhood.

% -----------------------------------------------------------------------
\section{Hyperparameters}

\begin{table}[h]
\centering
\caption{Tunable hyperparameters.}
\label{tab:hyperparams}
\begin{tabular}{lll}
\toprule
Parameter & Role & Starting value \\
\midrule
$W$       & Lookahead window
          & $150$--$300$ steps ($12.5$--$25$\,hr at 5-min resolution) \\
$\lambda$ & Reward-B global-term weight
          & $0.1$--$0.3$ \\
\bottomrule
\end{tabular}
\end{table}

\noindent
$W$ should be set to at least the median VM lifetime in the trace.
The future-knowledge ablation \cite{futurek} shows that the
marginal value of lookahead is low below $\approx\!150$ steps and
rises steeply above that threshold, justifying the lower bound.
$\lambda$ is only used in Reward B.  It should be kept small enough
that the local term still dominates the reward gradient across a
training batch; verify empirically before scaling to 96-host pods.

% -----------------------------------------------------------------------
\section{Potential Issues to Validate}

This section records the main formulation risks to check during
implementation and ablation.  These are not claims that the design
cannot work; they are the most likely places where the reward/state
proposal could fail to beat greedy or could produce misleading
training signal.

\subsection{Highest-Priority Risks}

\begin{enumerate}
    \item \textbf{Reward B's $\lambda$ term is only a coarse ablation.}
    The Reward B term
    $\max_{j \notin N(i)} \hat{c}_j(t)$ is not reconstructed exactly by
    the shared state.  This is acceptable for an ablation, but Reward B
    should not be described as a perfectly aligned state--reward design.

    \item \textbf{Reward B's $\lambda$ term may be action-independent at each step.}
    For a given VM arrival, the agent cannot change the loads of
    unreachable MPDs.  As written, the second reward term may therefore
    be constant with respect to the current action, which can add critic
    noise without helping the policy distinguish better placements.
    This should be checked empirically before relying on it as useful
    shaping.

    \item \textbf{Missing lifetime signal for the arriving VM.}
    The redesign distinguishes draining load from sticky load for
    already-placed VMs, but the state includes only the arriving VM size
    $x$, not the arriving VM's predicted lifetime.  If the new VM itself
    can be short-lived or long-lived, then the policy is missing a key
    variable needed to reason about whether the new placement will
    create future pressure.

    \item \textbf{No explicit feasibility / overflow term is documented.}
    The reward penalises projected load but does not explicitly encode a
    hard penalty for exceeding MPD capacity.  If infeasible allocations
    are blocked by an action mask or safety layer, that should be stated
    clearly.  If not, the policy may optimise the smooth peak surrogate
    while still learning unsafe behavior near capacity.
\end{enumerate}

\subsection{Secondary Risks}

\begin{enumerate}
    \item \textbf{Reward ranges should be reported per variant.}
    Reward A and Reward B have different lower bounds.  The text should
    keep those bounds separate rather than describing a single shared
    range for all reward variants.

    \item \textbf{Normalising by $D_\mathrm{pod}$ may weaken local signal.}
    The systems objective is driven by per-MPD peaks, but
    $D_\mathrm{pod}$ grows with pod size.  The same bad local placement
    can therefore produce a smaller reward difference in larger pods.
    A per-MPD capacity scale, or a fair-share normalization, may align
    better with the actual control objective.

    \item \textbf{The reward may remain close to a future-aware greedy rule.}
    The first term still behaves like ``minimise the post-placement
    projected local maximum over reachable MPDs.''  That is a reasonable
    heuristic, but it is not obviously the same as minimising the final
    trace-wide peak.  The formulation should be validated by checking
    whether greedy, random, and stronger baselines rank correctly under
    this reward on held-out traces.

    \item \textbf{Neighborhood contention features need empirical validation.}
    The new $P_j$ and $Q_j$ signals are intentionally simple summaries of
    local competition.  They may help beat greedy, but they should still
    be ablated to confirm that they improve decisions rather than just
    adding noise.

    \item \textbf{Slot ordering may create a permutation problem.}
    The observation is written as slot-indexed features
    $\phi_k^{(1:4)}$.  If MPDs are not placed into those slots using a
    canonical order, then equivalent neighborhoods may appear as
    different input vectors across episodes, making the MLP learn a
    harder permutation-sensitive task than intended.
\end{enumerate}

\subsection{Clarifying Questions to Resolve}

\begin{itemize}
    \item Is $D_\mathrm{pod}$ total pod DRAM, total MPD capacity, or a
    fair-share normalization constant?
    \item Is the arriving VM's lifetime known at decision time, or does
    deployment require a predictor for that quantity as well?
    \item Are capacity violations prevented by a hard mechanism outside
    the reward, or should the reward itself teach the policy that
    infeasible actions are unacceptable?
\end{itemize}

% -----------------------------------------------------------------------
\bibliographystyle{plain}
\begin{thebibliography}{9}
\bibitem{pond2023}
Huaicheng Li, Daniel S.\ Berger, Lisa Hsu, et al.
\newblock Pond: {CXL}-Based Memory Pooling Systems for Cloud Platforms.
\newblock In \textit{ASPLOS}, 2023.

\bibitem{futurek}
Richard Uzeel.
\newblock Perfect Future Knowledge.
\newblock Internal project report, 2026.
\end{thebibliography}

\end{document}